\section{Related Work}\label{sec:related}
\textbf{AutoML systems and benchmarks:} AutoML automates the choice of a learning algorithm and its hyperparameters, a task known as the CASH problem. Auto-WEKA treated algorithm choice and hyperparameter setting as one Bayesian optimization problem~\cite{thornton2013}, and TPOT searches over complete pipelines with genetic programming~\cite{olson2016}. Auto-sklearn warm-starts its search with meta-learning and combines the best models into an ensemble~\cite{feurer2015}, and its second version aims at use without manual settings~\cite{feurer2020}. H2O AutoML trains tree ensembles, linear models, and neural networks and stacks them~\cite{ledell2020}. AutoGluon-Tabular relies more on multi-layer stacking than on search~\cite{erickson2020}, and FLAML emphasizes a lightweight search~\cite{wang2019flaml}. Google Vizier is a black-box optimizer used for tuning, and it is not an end-to-end AutoML system~\cite{golovin2017}, and a survey covers the wider field~\cite{he2021}. Comparing such systems is difficult, because their search spaces, objectives, and budgets differ. Gijsbers et al. built an open benchmark with 71 classification and 33 regression tasks and ran nine frameworks on it~\cite{gijsbers2022}. Average ranks were mostly not significantly different, and rankings depended on data properties such as size and class balance. The tasks are general tabular problems and contain no IoT or security data.

\textbf{Automated model selection for intrusion and IoT detection:} Hand-built models remain common in intrusion detection. Deep networks, boosting, and ensembles reach high scores on benchmark data~\cite{vinayakumar2019,tama2019,bibers2024}. Yang and Shami propose an AutoML-based intrusion detection system for next-generation networks, with automated feature engineering and Bayesian optimization~\cite{yang2024automl}. Injadat et al. chain oversampling, feature selection, and hyperparameter optimization in one pipeline and test it on CICIDS2017 and UNSW-NB15~\cite{injadat2021}. AID4I applies AutoML to industrial IoT intrusion detection and automates preprocessing, feature selection, model selection, and tuning~\cite{sezgin2023}. AutoAI-IDS trains 64 configurations from seven model families and uses meta-learning to select a model for each sample on NSL-KDD and CICIDS-2017, and it benchmarks H2O AutoML and Auto-sklearn as references~\cite{guntur2025autoai}. In our earlier study, we compared the same two frameworks on the same two datasets~\cite{velnambi2026}. The weighted F1 score of H2O AutoML reached 0.99 on CICIDS-2017, but its macro-F1 was 0.68, and the two frameworks saw different training data, so that their gap could not be assigned to the frameworks. These studies address network traffic, and none of them compares AutoML frameworks on IoT sensor data to our knowledge.

\textbf{IoT anomaly detection and explainability:} IoT anomaly detection covers energy use in buildings~\cite{himeur2021}, agriculture and manufacturing~\cite{abdallah2021agri,abdallah2023sensors}, imbalanced smart-manufacturing data~\cite{fathy2021}, and anomalous machine sounds~\cite{koizumi2020}. For network traffic, N-BaIoT provides labeled botnet attacks from real devices~\cite{meidan2018nbaiot}, and Kitsune detects attacks online with an ensemble of autoencoders~\cite{mirsky2018kitsune}. Explainable AI shows which features drive the decision of a model. XAI-IoT compares ten models and seven explanation methods on MEMS and N-BaIoT, and it finds that the best model depends on the dataset~\cite{gummadi2024xaiiot}. Related studies extend the analysis to white-box explanation methods~\cite{gummadi2025whitebox} and to intrusion detection~\cite{arreche2024xaiids}. XAI-IoT uses hand-built models and a 70/30 split and does not audit the data. We reuse its two datasets, so that our results can be compared with theirs.

\textbf{Data quality and evaluation pitfalls:} Benchmark flaws can change conclusions. Tavallaee et al. showed that KDD CUP 99 contains problems that distort evaluation, which led to NSL-KDD~\cite{tavallaee2009}. Engelen et al. found errors in traffic generation, flow extraction, and labeling in CICIDS2017 and corrected more than 20\% of the traces~\cite{engelen2021}. Lanvin et al. later reported packet misorder, packet duplication, and attacks that were not labeled~\cite{lanvin2023}. Arp et al. surveyed 30 security papers and identified pitfalls such as sampling bias and unrealistic evaluation settings~\cite{arp2022}. Kapoor and Narayanan document data leakage in machine-learning-based science~\cite{kapoor2023}, and Northcutt et al. show that label errors in test sets can change model rankings~\cite{northcutt2021}. For dependent data, blocked cross-validation can give more realistic estimates than random folds~\cite{roberts2017}.

Our work differs from these works in that it combines the three lines above. It applies a controlled comparison of AutoML frameworks to IoT data, it audits the released datasets before it ranks methods, and it reports cost and explanations next to detection quality. The audit of the released datasets and the effect of the audit on the ranking of the frameworks are, to our knowledge, not covered by prior work.
